% ---- pg-382 （书页 370） ----

\noindent\hfill 续表

\begin{mvtable}
\begin{tabular}{lcccccccc}
\toprule
国家和地区 & 100m & 200m & 400m & 800m & 1 500m & 5 000m & 10 000m & 马拉松 \\
 & (s) & (s) & (s) & (min) & (min) & (min) & (min) & (min) \\
\midrule
卢森堡 & 10.41 & 20.77 & 47.90 & 1.76 & 3.67 & 13.64 & 28.77 & 134.03 \\
马来西亚 & 10.30 & 20.92 & 46.41 & 1.79 & 3.76 & 14.11 & 29.50 & 149.27 \\
毛里求斯 & 10.13 & 20.06 & 44.69 & 1.80 & 3.83 & 14.15 & 29.84 & 143.07 \\
墨西哥 & 10.21 & 20.40 & 44.31 & 1.78 & 3.63 & 13.13 & 27.14 & 127.19 \\
缅甸 & 10.64 & 21.52 & 48.63 & 1.80 & 3.80 & 14.19 & 29.62 & 139.57 \\
荷兰 & 10.19 & 20.19 & 45.68 & 1.73 & 3.55 & 13.22 & 27.44 & 128.31 \\
新西兰 & 10.11 & 20.42 & 46.09 & 1.74 & 3.54 & 13.21 & 27.70 & 128.59 \\
挪威 & 10.08 & 20.17 & 46.11 & 1.71 & 3.62 & 13.11 & 27.54 & 130.17 \\
巴布亚新几内亚 & 10.40 & 21.18 & 46.77 & 1.80 & 4.00 & 14.72 & 31.36 & 148.13 \\
菲律宾 & 10.57 & 21.43 & 45.57 & 1.80 & 3.82 & 13.97 & 29.04 & 138.44 \\
波兰 & 10.00 & 19.98 & 44.62 & 1.72 & 3.59 & 13.29 & 27.89 & 129.23 \\
葡萄牙 & 9.86 & 20.12 & 46.11 & 1.75 & 3.50 & 13.05 & 27.21 & 126.36 \\
罗马尼亚 & 10.21 & 20.75 & 45.77 & 1.76 & 3.57 & 13.25 & 27.67 & 132.30 \\
俄罗斯 & 10.11 & 20.23 & 44.60 & 1.71 & 3.54 & 13.20 & 27.90 & 129.16 \\
萨摩亚 & 10.78 & 21.86 & 49.98 & 1.94 & 4.01 & 16.28 & 34.71 & 161.50 \\
新加坡 & 10.37 & 21.14 & 47.60 & 1.84 & 3.86 & 14.96 & 31.32 & 144.22 \\
西班牙 & 10.17 & 20.59 & 44.96 & 1.73 & 3.48 & 13.04 & 27.24 & 127.23 \\
瑞典 & 10.18 & 20.43 & 45.54 & 1.76 & 3.61 & 13.29 & 27.93 & 130.38 \\
瑞士 & 10.16 & 20.41 & 44.99 & 1.71 & 3.53 & 13.13 & 27.90 & 129.56 \\
中国台湾 & 10.36 & 20.81 & 46.72 & 1.79 & 3.77 & 13.91 & 29.20 & 134.35 \\
泰国 & 10.23 & 20.69 & 46.05 & 1.81 & 3.77 & 14.25 & 29.67 & 139.33 \\
土耳其 & 10.38 & 21.04 & 46.63 & 1.78 & 3.59 & 13.45 & 28.33 & 130.25 \\
美国 & 9.78 & 19.32 & 43.18 & 1.71 & 3.46 & 12.97 & 27.21 & 125.38 \\
\bottomrule
\end{tabular}
\end{mvtable}

\vspace{0.5em}
\noindent{\small 资料来源: \textit{IAAF/ATES Track and Field Statistics Handbook} for the Helsinki 2005 Olympics. Courtesy of Ottavio Castellini.}

\noindent 8.21\quad 参见练习 8.20, 将\hyperref[tab:8.6]{表 8.6} 中的国家和地区男子径赛记录的速度换成用米/秒表示. 注意: 记录中 800 米、1 500 米、5 000 米、10 000 米和马拉松是以分钟计的. 其中, 马拉松是 26.2 英里或 42 195 米长. 用速度的协方差矩阵 $\mathbf{S}$ 作主成分分析, 同练习 8.20 的结果比较一下, 你认为哪一个分析得更好? 为什么?

\noindent 8.22\quad 考虑\hyperref[tab:1.10]{表 1.10} 的公牛数据, 利用这 7 个变量 YrHgt, FtFtBody, PrctFFB, Frame, BkFat, SaleHt 及 SaleWt, 由协方差矩阵 $\mathbf{S}$ 和相关系数矩阵 $\mathbf{R}$ 作主成分分析, 你的分析中应包括以下各项:

(a) 确定能有效地综合样本变异性的恰当的成分个数. 构造崖底碎石图帮助你的求解.

(b) 解释样本主成分.

(c) 你认为由以上 7 个变量的数据可能得到一个“身体尺寸”或“身体外形”指标吗? 请说明一下.

(d) 用前两个主成分的值, 将这些数据画在 $\hat{y}_1$ 表示竖轴, $\hat{y}_2$ 表示横轴的二维空间上, 你能区分代表三种牛的组吗? 有没有离群值?
